\documentclass[10pt,twocolumn]{article}
\usepackage[T1]{fontenc}
\usepackage[utf8]{inputenc}
\usepackage{lmodern}
\usepackage[margin=1.8cm,top=2.4cm,bottom=2cm,columnsep=0.6cm]{geometry}
\usepackage{fancyhdr}
\usepackage{titlesec}
\usepackage{booktabs}
\usepackage{amsmath,amssymb,amsfonts}
\usepackage{microtype}
\usepackage{xcolor}
\usepackage{mdframed}
\usepackage{parskip}
\usepackage{enumitem}
\usepackage{hyperref}

\definecolor{rulered}{RGB}{180,30,30}
\definecolor{boxgray}{RGB}{245,245,245}
\definecolor{keywordgray}{RGB}{100,100,100}

\pagestyle{fancy}
\fancyhf{}
\fancyhead[L]{\textsc{\small Claude Mag}}
\fancyhead[C]{\small\textit{Machine Learning Research Digest}}
\fancyhead[R]{\small 2026-10-12}
\fancyfoot[C]{\small\thepage}
\renewcommand{\headrulewidth}{0.8pt}

\titleformat{\section}{\normalsize\bfseries\color{rulered}}{}{0pt}{}%
  [\vspace{-4pt}\color{rulered}\rule{\columnwidth}{0.4pt}]
\titlespacing{\section}{0pt}{8pt}{4pt}

\hypersetup{colorlinks=true,urlcolor=rulered,linkcolor=black}

\begin{document}
\setlength{\parindent}{0pt}
\setlength{\parskip}{4pt}

{\small\color{keywordgray}\textbf{Keywords:}
  multimodal learning \textbullet{}
  tensor decomposition \textbullet{}
  synergy \textbullet{}
  contrastive learning}\par\vspace{4pt}

{\LARGE\bfseries\raggedright
  HRIL: Learning Multimodal Synergy via Higher-Order Tensor Modeling}\par\vspace{4pt}

{\large\itshape\raggedright
  Tucker Decomposition of a Cross-Moment Tensor Captures the Emergent
  Information That Pairwise Contrastive Objectives Miss}\par\vspace{6pt}

{\small\textbf{By} Qun Dai, Zhao Kang et al.}\par\vspace{2pt}

{\small\color{keywordgray}arXiv:2610.12393 \textbullet{} NeurIPS 2026 \textbullet{} October 2026}
\par\vspace{2pt}\noindent{\color{rulered}\rule{\columnwidth}{0.6pt}}\vspace{6pt}

Multimodal contrastive learning fails to capture \emph{synergistic} information---the
signal that is only present in the joint combination of three or more
modalities. HRIL introduces a cross-moment tensor over modality embeddings
and factorises it with Tucker decomposition inside a contrastive objective,
with a regulariser that prevents the core tensor from collapsing to
pairwise terms.

\begin{mdframed}[backgroundcolor=boxgray,linecolor=rulered,linewidth=1pt,
    innerleftmargin=6pt,innerrightmargin=6pt,
    innertopmargin=6pt,innerbottommargin=6pt]
\textbf{\small AT A GLANCE}\par\vspace{4pt}
\begin{description}[leftmargin=0pt,labelsep=4pt,itemsep=2pt,parsep=0pt,
    font=\small\bfseries]
  \item[Problem:] Pairwise contrastive losses miss synergistic information
    that only exists in the joint combination of three or more modalities.
  \item[Why it matters:] Many real-world multimodal tasks rely on
    cross-modal synergy; ignoring it leaves performance on the table.
  \item[Method:] Outer-product cross-moment tensor over modality
    embeddings, compressed by Tucker decomposition, with a diagonal-suppressing
    regulariser.
  \item[Result:] Consistent gains over multimodal contrastive baselines
    on synergy-dominated benchmarks; accepted to NeurIPS 2026.
  \item[Evidence:] Ablation confirms the synergy regulariser is necessary;
    Tucker rank controls the expressiveness--cost trade-off.
\end{description}
\end{mdframed}

\section{The Big Picture}

The information-theoretic decomposition of multimodal data identifies three
components: \emph{redundant} (shared across modalities), \emph{unique}
(present in one modality only), and \emph{synergistic} (present only in
the joint combination). Standard contrastive objectives such as InfoNCE
maximise pairwise mutual information, capturing redundant and unique
information but systematically missing synergy.

HRIL (Higher-order Representation Interaction Learning) addresses this gap
by constructing a tensor representation that explicitly encodes interactions
of all orders up to $M$ for $M$ modalities, then compressing it efficiently
via Tucker decomposition. The result is a contrastive framework that is
provably sensitive to synergistic effects.

\section{Why Existing Approaches Fall Short}

\textbf{Pairwise contrastive losses} (InfoNCE, SimCLR) consider only two
modalities at a time, so terms requiring all $M$ modalities simultaneously
are never optimised.

\textbf{Late fusion} combines independent unimodal embeddings at inference
time, encoding no interactions before the final decision.

\textbf{Bilinear tensor fusion} captures pairwise interactions and scales
quadratically with the number of modalities, missing higher-order terms
and becoming intractable for large $M$.

\section{Core Mechanism}

HRIL forms the outer product of modality embeddings:
\[
\mathcal{T}^{(i)} = \bigotimes_{m=1}^{M} \mathbf{z}_m^{(i)},
\quad \mathcal{T}^{(i)} \in \mathbb{R}^{d_1 \times \cdots \times d_M}
\]
and compresses it via Tucker decomposition into a core tensor
$\mathcal{G} \in \mathbb{R}^{r_1 \times \cdots \times r_M}$ with factor
matrices $\mathbf{U}_m \in \mathbb{R}^{d_m \times r_m}$. A regulariser
penalises the diagonal entries of $\mathcal{G}$ (which correspond to
unimodal and pairwise terms), forcing the model to allocate capacity to
genuinely multi-way interactions.

\section{Technical Formulation}

Let $f_m : \mathcal{X}_m \to \mathbb{R}^{d_m}$ be the encoder for modality
$m$, with output $\mathbf{z}_m^{(i)} = f_m(x_m^{(i)})$.

Tucker factorisation approximates the cross-moment tensor:
\[
\hat{\mathcal{T}} = \mathcal{G} \times_1 \mathbf{U}_1 \cdots \times_M \mathbf{U}_M,
\quad \|\mathcal{T} - \hat{\mathcal{T}}\|_F \leq \epsilon
\]

The synergy regulariser penalises diagonal concentration:
\[
\mathcal{L}_\text{syn} = \sum_{r} \mathcal{G}_{r,r,\ldots,r}^2
\]

The total objective balances contrastive alignment and synergy:
\[
\mathcal{L} = \mathcal{L}_\text{InfoNCE} + \lambda \mathcal{L}_\text{syn}
\]
where $\lambda > 0$ is a hyperparameter.

\section{Learning or Inference Procedure}

\textbf{Training:}
\begin{enumerate}[itemsep=2pt,parsep=0pt]
  \item Encode each modality: $\mathbf{z}_m^{(i)} = f_m(x_m^{(i)})$.
  \item Form the outer-product tensor $\mathcal{T}^{(i)}$.
  \item Project through Tucker factors to obtain the core $\mathcal{G}^{(i)}$.
  \item Evaluate the contrastive loss over positive/negative pairs.
  \item Add $\lambda\mathcal{L}_\text{syn}$ and backpropagate through
    Tucker factors and encoders jointly.
\end{enumerate}

\textbf{Inference:} The Tucker-compressed representation provides a
fused embedding for downstream classification or retrieval.

\section{What the Guarantee Says}

Tucker decomposition provides a controlled approximation: the Frobenius
reconstruction error is bounded by the residual from truncating the Tucker
rank. The synergy regulariser provably prevents the core from degenerating
to a diagonal form (which recovers bilinear fusion).

\section{Experimental Findings}

\begin{itemize}[itemsep=2pt,parsep=0pt]
  \item HRIL outperforms multimodal contrastive baselines on synergy-dominated
    benchmarks, with gains largest when no pair of modalities alone
    captures the label.
  \item Performance scales with Tucker rank $(r_1, \ldots, r_M)$ up to a
    saturation point, providing a tuneable compute--accuracy trade-off.
  \item Gains transfer across different encoder architectures
    and modality combinations.
\end{itemize}

\section{Ablations and Interpretation}

\textbf{Without synergy regulariser:} The core concentrates on diagonals,
recovering bilinear-fusion behaviour and losing most of the HRIL advantage.

\textbf{Tucker rank:} Lowering the rank reduces cost with minor accuracy
loss; higher ranks plateau beyond a problem-dependent threshold.

\textbf{Number of modalities:} The benefit over pairwise methods grows
with $M$, consistent with the theoretical argument that synergy is a
higher-order effect.

\section*{Reference}

Qun Dai, Zhao Kang et al.\
\textbf{HRIL: Learning Multimodal Synergy via Higher-Order Tensor Modeling}.
NeurIPS 2026. arXiv:2610.12393, October 2026.\\
\url{https://arxiv.org/abs/2610.12393}

\end{document}
